\documentclass[11pt,a4paper]{article}

\usepackage[utf8]{inputenc}
\usepackage[T1]{fontenc}
\usepackage[french]{babel}

\usepackage{geometry}
\geometry{
  a4paper,
  left=25mm,
  right=25mm,
  top=25mm,
  bottom=25mm
}

\usepackage{xcolor}
\usepackage{titlesec}
\usepackage{fancyhdr}
\usepackage{listings}
\usepackage{hyperref}
\usepackage{sectsty}

% Couleurs
\definecolor{primary}{RGB}{43, 76, 126}
\definecolor{secondary}{RGB}{86, 110, 157}
\definecolor{codebg}{RGB}{245, 245, 245}
\definecolor{codeborder}{RGB}{220, 220, 220}

% Configuration des hyperliens
\hypersetup{
    colorlinks=true,
    linkcolor=primary,
    filecolor=magenta,      
    urlcolor=primary,
    pdftitle={Compte Rendu TP RDF/OWL},
}

% Configuration des titres
\titleformat{\section}{\Large\bfseries\color{primary}}{}{0em}{}[\titlerule]
\titleformat{\subsection}{\large\bfseries\color{secondary}}{\thesubsection}{1em}{}

% Configuration de l'entête et pied de page
\pagestyle{fancy}
\fancyhf{}
\fancyhead[L]{\textcolor{gray}{TP Intégration de données - RDF/OWL}}
\fancyhead[R]{\textcolor{gray}{AMMARI Reda \& ALJIBBAOUI Tala}}
\fancyfoot[C]{\thepage}
\renewcommand{\headrulewidth}{0.4pt}
\renewcommand{\headrule}{\hbox to\headwidth{\color{gray}\leaders\hrule height \headrulewidth\hfill}}

% Configuration Listings pour le code Turtle
\lstset{
    backgroundcolor=\color{codebg},
    basicstyle=\ttfamily\small,
    breaklines=true,
    captionpos=b,
    commentstyle=\color{green!60!black},
    keywordstyle=\color{blue},
    stringstyle=\color{red},
    frame=single,
    rulecolor=\color{codeborder},
    xleftmargin=10pt,
    xrightmargin=10pt,
    tabsize=2,
    aboveskip=15pt,
    belowskip=15pt
}

\begin{document}

\begin{titlepage}
    \centering
    \vspace*{4cm}
    
    \rule{\textwidth}{1pt}\vspace{0.5cm}
    {\huge \bfseries \textcolor{primary}{Compte Rendu de TP}} \\ \vspace{0.3cm}
    {\Large \bfseries Modélisation et Intégration de données : RDF et OWL}
    \vspace{0.5cm}\rule{\textwidth}{1pt}
    
    \vspace{2cm}
    
    {\Large \textbf{Auteurs :}} \\ \vspace{0.2cm}
    {\large AMMARI Reda \\ ALJIBBAOUI Tala}
    
    \vspace{2cm}
    
    {\Large \textbf{Module :}} \\ \vspace{0.2cm}
    {\large Traitement Sémantique des Données}
    
    \vfill
    {\large Année Universitaire 2025/2026}
    
\end{titlepage}

\newpage

\section*{Exercice 2 : Focus sur les Connaissances : RDF-Schema et OWL}

\subsection*{1. Identification des vocabulaires}

\begin{itemize}
    \item \textbf{rdf} : utilisé pour la structure de base (définir le type des objets).
    \item \textbf{rdfs} : utilisé pour le schéma (créer des hiérarchies de classes et des règles sur les propriétés).
    \item \textbf{owl} : utilisé pour les ontologies complexes (définir précisément la nature des propriétés).
    \item \textbf{movies} : vocabulaire local spécifique au domaine du cinéma (films, acteurs, réalisateurs).
    \item \textbf{dbp} : vocabulaire externe (DBpedia) utilisé pour identifier des personnes réelles comme Alfred Hitchcock.
\end{itemize}

\subsection*{2. Séparation des triplets}

\textbf{Connaissances Ontologiques (Le Schéma) :}
Ces triplets définissent les classes, les hiérarchies et les contraintes sur les propriétés (domaines et portées) que nous avons placés dans le fichier \texttt{owlExo2Schema.ttl}.

\begin{lstlisting}[language=SQL]
@prefix movies: <http://www.lirmm.fr/ulliana/movies#> .
@prefix rdfs: <http://www.w3.org/2000/01/rdf-schema#> .
@prefix rdf: <http://www.w3.org/1999/02/22-rdf-syntax-ns#> .
@prefix owl: <http://www.w3.org/2002/07/owl#> .

movies:directedBy rdfs:domain movies:Movie .
movies:title rdfs:domain movies:Movie .
movies:playsIn rdfs:domain movies:Actor .
movies:playsIn rdfs:range movies:Movie .
movies:Actor rdfs:subClassOf movies:Artist .
movies:Director rdfs:subClassOf movies:Artist .
movies:title rdf:type owl:DatatypeProperty .
movies:directedBy rdfs:range movies:Director .
\end{lstlisting}

\newpage

\textbf{Triplets représentant des données (Les Faits) :}
Ces triplets décrivent des individus spécifiques (m1, m2, m3, a1) et leurs relations concrètes que nous avons placés dans le fichier \texttt{owlExo2Data.ttl}.

\begin{lstlisting}[language=SQL]
@prefix movies: <http://www.lirmm.fr/ulliana/movies#> .
@prefix rdf: <http://www.w3.org/1999/02/22-rdf-syntax-ns#> .
@prefix dbp: <http://dbpedia.org/> .

movies:m3 movies:directedBy dbp:Alfred_Hitchcock .
movies:m2 movies:title "Vertigo" .
movies:m1 rdf:type movies:Movie .
movies:a1 movies:playsIn movies:m4 .
\end{lstlisting}

\subsection*{3. Triplets inférés}

En combinant le schéma (règles) et les données (faits), nous obtenons les nouvelles informations suivantes :

\textbf{Inférences basées sur les Domaines (\texttt{rdfs:domain})}
\begin{itemize}
    \item \texttt{movies:m2 rdf:type movies:Movie .} (Car m2 possède un \texttt{title})
    \item \texttt{movies:m3 rdf:type movies:Movie .} (Car m3 est sujet de \texttt{directedBy})
    \item \texttt{movies:a1 rdf:type movies:Actor .} (Car a1 est sujet de \texttt{playsIn})
\end{itemize}

\textbf{Inférences basées sur les Portées (\texttt{rdfs:range})}
\begin{itemize}
    \item \texttt{dbp:Alfred\_Hitchcock rdf:type movies:Director .} (Car il est l'objet de \texttt{directedBy})
    \item \texttt{movies:m4 rdf:type movies:Movie .} (Car il est l'objet de \texttt{playsIn})
\end{itemize}

\textbf{Inférences basées sur la Hiérarchie (\texttt{rdfs:subClassOf})}
\begin{itemize}
    \item \texttt{dbp:Alfred\_Hitchcock rdf:type movies:Artist .} (Car c'est un Director et tout réalisateur est un artiste)
    \item \texttt{movies:a1 rdf:type movies:Artist .} (Car c'est un Actor et tout acteur est un artiste)
\end{itemize}

\subsection*{4. Vérification}
Oui, en chargeant les données avec Jena, le modèle contient bien toutes les inférences listées ci-dessus.

\newpage

\section*{Exercice 3 : De RDF-Schema à OWL (Extensions du modèle)}

\subsection*{1. Propriété Symétrique}
\textbf{Propriété choisie :} \texttt{:isFriendOf} \\
Dans \texttt{owlDemoSchema.ttl}, nous avons défini \texttt{:isFriendOf} comme de type \texttt{owl:SymmetricProperty}.
Dans \texttt{owlDemoData.ttl}, nous avons ajouté le fait : \texttt{:Alice :isFriendOf :Bob}.

\textbf{Inférence constatée :} \\
En exécutant le raisonneur OWL (\texttt{RDFOWLReasoning.java}), le système déduit automatiquement que la relation inverse est vraie, générant le triplet inféré suivant :
\begin{quote}
\texttt{http://example.org/Bob http://example.org/isFriendOf http://example.org/Alice}
\end{quote}

\subsection*{2. Propriété Fonctionnelle}
\textbf{Propriété choisie :} \texttt{:hasPrimaryDoctor} \\
Dans \texttt{owlDemoSchema.ttl}, nous avons défini \texttt{:hasPrimaryDoctor} comme de type \texttt{owl:FunctionalProperty}.
Dans \texttt{owlDemoData.ttl}, nous avons déclaré que le patient \texttt{:Alice\_Rel} a pour médecin principal \texttt{:DrSmith} et \texttt{:DrJohnSmith}.

\textbf{Inférence constatée :} \\
En exécutant le raisonneur OWL, le système déduit que puisqu'un patient ne peut avoir qu'un seul médecin principal (propriété fonctionnelle), les deux ressources \texttt{:DrSmith} et \texttt{:DrJohnSmith} doivent en réalité représenter la même entité. Il génère les triplets d'équivalence suivants :
\begin{quote}
\texttt{http://example.org/DrSmith owl:sameAs http://example.org/DrJohnSmith}\\
\texttt{http://example.org/DrJohnSmith owl:sameAs http://example.org/DrSmith}
\end{quote}

\subsection*{3. Propriété Inverse}
\textbf{Propriétés choisies :} \texttt{:isTreatedBy} comme inverse de \texttt{:treats}. \\
Dans \texttt{owlDemoSchema.ttl}, nous avons défini \texttt{:isTreatedBy} avec \texttt{owl:inverseOf :treats}.
Dans \texttt{owlDemoData.ttl}, nous avons déclaré le fait : \texttt{:Charlie\_Rel :isTreatedBy :DrHouse}.

\textbf{Inférence constatée :} \\
En exécutant le raisonneur OWL, le système déduit automatiquement la relation dans l'autre sens. Il génère le triplet inverse suivant :
\begin{quote}
\texttt{http://example.org/DrHouse http://example.org/treats http://example.org/Charlie\_Rel}
\end{quote}

\subsection*{4. Propriété Transitive}
\textbf{Propriété choisie :} \texttt{:isSuperiorOf} \\
Dans \texttt{owlDemoSchema.ttl}, nous avons défini \texttt{:isSuperiorOf} comme \texttt{owl:TransitiveProperty}.
Dans \texttt{owlDemoData.ttl}, nous avons créé une chaîne hiérarchique avec deux faits : \texttt{:AliceBoss :isSuperiorOf :BobManager} et \texttt{:BobManager :isSuperiorOf :CharlieWorker}.

\textbf{Inférence constatée :} \\
En exécutant le raisonneur OWL, le système déduit par transitivité que la relation s'étend du premier au dernier maillon de la chaîne. Il génère le triplet attendu :
\begin{quote}
\texttt{http://example.org/AliceBoss http://example.org/isSuperiorOf \\ http://example.org/CharlieWorker}
\end{quote}

\end{document}
